\chapter{Methodology}

\section{Overview and Experimental Design}

This chapter describes how the benchmark was constructed and how every reported
number was produced.
The goal is a fair and reproducible comparison of Ascon-AEAD128 and AES-128-GCM
under matched conditions on two constrained microcontrollers.
Fairness here has a precise meaning: both algorithms run as software-only
implementations, both are measured through an identical timing harness, both are
compiled with the same toolchain and the same set of optimization levels, and the
one-time setup cost is excluded from the timed region for both in the same way.

The experimental design varies four factors: the algorithm (Ascon-AEAD128 or
AES-128-GCM), the hardware platform (Cortex-M0 or Cortex-M4F), the input
message size (swept across several orders of magnitude), and the compiler
optimization level (swept from no optimization through size optimization).
Each combination is measured many times and reduced to summary statistics
offline.
Code size is measured separately, as a static property of the compiled binaries
rather than a runtime measurement.

\section{Target Platforms}

Two STMicroelectronics Nucleo-64 development boards represent the low and
middle of the constrained range.
Their relevant properties are summarized in Table~\ref{tbl:platforms}.
Both boards expose an on-board ST-LINK debugger with a USB virtual COM port
routed to USART2, which carries the measurement output.

\begin{table}[htbp]
  \centering
  \caption{Target platform properties}
  \label{tbl:platforms}
  \begin{tabular}{lll}
    \toprule
    Property              & NUCLEO-F072RB     & NUCLEO-L476RG    \\
    \midrule
    Microcontroller       & STM32F072RBT6     & STM32L476RGT6    \\
    Core                  & Cortex-M0         & Cortex-M4F       \\
    Maximum clock         & 48\,MHz           & 80\,MHz          \\
    Flash                 & 128\,KB           & 1\,MB            \\
    SRAM                  & 16\,KB            & 128\,KB          \\
    Hardware AES          & None              & None             \\
    Cycle counter (DWT)   & Absent            & Present          \\
    Instruction prefetch  & None              & ART accelerator  \\
    Floating-point unit   & None              & Single precision \\
    \bottomrule
  \end{tabular}
\end{table}

\subsection{NUCLEO-F072RB (Cortex-M0)}

The F072 carries a Cortex-M0 core with no hardware multiplier for wide
operands, no instruction cache, no branch predictor, and no floating-point
unit.
It represents the lower end of the device range, where every cycle and every
kilobyte is scarce.
Its 16\,KB of SRAM is the binding constraint on this study, and it limits the
largest message size that can be benchmarked on this board.

\subsection{NUCLEO-L476RG (Cortex-M4F)}

The L476 carries a Cortex-M4F core with a single-cycle hardware multiplier, a
single-precision floating-point unit, and the ST ART accelerator, which
prefetches and caches instructions from Flash.
It represents the middle of the range.
The L476 has substantially more memory than the F072, so memory is not a
constraint on this board, but it is included to expose how the relative
performance of the two algorithms depends on the processor architecture.

The specific part on the NUCLEO-L476RG does not contain a hardware AES block;
only the L486 and related parts in the same family include one.
This absence is useful for the study, because it means there is no hardware
cryptographic path that could be engaged accidentally, so the software
comparison is consistent on both boards.

\subsection{Clocking and the Matched-Frequency Decision}

Both boards run at 48\,MHz for the timing comparison.
The L476 is capable of 80\,MHz, so it is intentionally underclocked to the
F072 ceiling.
This choice removes clock frequency as a variable: any difference in measured
time between the two boards at the same clock then reflects the processor
architecture alone, not a difference in operating frequency.

The measurements are reported primarily as cycle counts rather than wall-clock
time, since the timer increments once per processor clock (see
Section~\ref{sec:timing-source}), so one timer tick equals one clock cycle and
the cycle count is independent of the operating frequency.

\section{Toolchain and Build System}

All firmware was built with a single toolchain to keep the comparison
consistent.
The compiler is \texttt{arm-none-eabi-gcc} version 14.3.1, supplied with the
STM32CubeCLT bundle.
The build is driven by CMake with the Ninja generator, both bundled with
CubeCLT.
Project skeletons, including peripheral initialization and the linker script,
were generated with STM32CubeMX configured for CMake output.
The host operating system was Windows.

The base compiler flags follow the CubeMX template.
They include \texttt{-ffunction-sections}, \texttt{-fdata-sections}, and linker
garbage collection (\texttt{-Wl,--gc-sections}), the size-oriented newlib-nano
C library (\texttt{--specs=nano.specs}), and the architecture flags appropriate
to each core.
On the L476 these include the hardware floating-point ABI, although
floating-point arithmetic is not used by either algorithm.

The optimization level is the one build variable that changes across the
measurement matrix.
It is controlled through a single CMake cache variable set per build preset.
Each preset reconfigures the build at one optimization level, from \texttt{-O0}
through \texttt{-Os}, so that the same source produces a family of binaries
that differ only in optimization.

\section{Cryptographic Implementations}

Both implementations are taken from existing open-source projects rather than
written for this study.
This keeps the comparison representative of what a developer would actually
deploy, and it avoids the risk of an unrepresentative hand-written
implementation favoring one side.

\subsection{AES-128-GCM (mbedTLS)}

AES-128-GCM is provided by mbedTLS version 3.6.x~\cite{mbedTLS}.
mbedTLS was selected because it is a widely deployed, well-maintained library
with a complete and verified AES-GCM implementation, which lends the comparison
credibility.
The library was reduced to a minimal configuration containing only the modules
required for AES-128-GCM.

Three configuration choices affect the measurements and are disclosed here,
because the reported numbers are only interpretable alongside them.

First, the small four-bit GHASH multiplication table is used.
The larger eight-bit table is faster per byte but enlarges the GCM context by
roughly 3.8\,KB of RAM.
The small table was chosen because it represents a genuinely lightweight
configuration of AES-GCM, matching the constrained-device framing of this work.
As a consequence, the reported AES-GCM throughput represents the memory-frugal
end of what mbedTLS can achieve rather than its speed ceiling.
This point is revisited in the discussion.

Second, the AES lookup tables are placed in Flash rather than generated in RAM
at initialization.
This shifts roughly 8\,KB from RAM to Flash, which is the realistic trade-off
on a part with only 16\,KB of SRAM.
It affects the code-size results but has no meaningful effect on runtime.

Third, only the 128-bit key length is compiled.
The 192-bit and 256-bit key schedules are removed.
Because only AES-128 is benchmarked, this has no effect on runtime and only a
small effect on Flash size.

\subsection{Ascon-AEAD128 (Reference Implementation)}

Ascon-AEAD128 is provided by the official ASCON reference implementation in
C~\cite{ASCONSoftware}.
The reference variant was chosen over the optimized variants for portability and
because it corresponds directly to the standardized specification.
The implementation uses a 16-byte key, a 16-byte nonce, and produces a 16-byte
authentication tag, with a rate of 16 bytes.

The choice of the reference variant has a visible effect on code size at high
optimization levels.
This is reported and discussed rather than hidden, because it reflects what a
developer integrating the standard reference code would observe.

\subsection{Uniform Wrapper Interface}

Each implementation is reached through a thin wrapper with a matching
signature, so that the benchmark harness calls both algorithms the same way.
The AES wrapper separates key setup from encryption, so that the one-time key
schedule can be hoisted out of the timed region.
The ASCON single-pass AEAD interface has no separable key schedule, so its
encryption call is timed in full.
Neither wrapper performs additional copying inside the timed path beyond what
the underlying library requires.

\section{Correctness Verification}

No timing measurement was trusted until the implementation producing it was
shown to be correct.
Each algorithm was checked against a published Known Answer Test vector, the
NIST GCM vector for AES-128-GCM and a Lightweight Cryptography Known Answer Test
vector for Ascon-AEAD128, confirming both the ciphertext and the authentication
tag.
Each was additionally checked with a round-trip test, where decryption returns
the original plaintext, and a forgery-rejection test, where a modified tag is
correctly rejected.
Both verification routines run on the target hardware at startup, before the
benchmark begins, print a pass result over the serial link, and halt on failure
so that no subsequent timing for that build is interpreted.

\section{Measurement Methodology}
\label{sec:timing-source}

Both boards use the TIM2 timer as the timing source.
TIM2 is a 32-bit timer on both families, configured identically on each board as
a free-running up-counter driven by the internal clock with the prescaler set to
zero, so the counter increments once per processor clock cycle and one tick
equals one cycle.
At 48\,MHz the 32-bit counter wraps approximately every 89 seconds, far longer
than any single measurement, and all elapsed-time computations use unsigned
32-bit subtraction, which handles a single wrap during a measurement correctly.
The Cortex-M0 on the F072 does not provide the Data Watchpoint and Trace cycle
counter available on the Cortex-M4F, so using TIM2 on both boards gives a
comparison between identically configured peripherals rather than between two
architecturally different timing mechanisms.

The benchmark macro disables all maskable interrupts for the duration of the
measured region and re-enables them afterward, which prevents the periodic
system tick interrupt and any peripheral interrupts from polluting short
measurements.
The effect was confirmed during validation, where masking reduced the
per-iteration variation from roughly 28 ticks to zero on a reference workload.
Each measurement also runs the workload once before the timed iterations begin
and discards the result.
This warmup primes the L476 ART accelerator and any related microarchitectural
state, on which the first run differs from the steady state by roughly 1.5\% at
small sizes.
On the F072 there is no measurable warmup effect, because the Cortex-M0 has no
instruction cache or branch predictor that could be cold.

At higher optimization levels the compiler can remove a computation whose result
is never used.
To prevent this, each workload writes one byte of its output into a
\texttt{volatile} variable, which forces the compiler to produce the output
because it cannot prove the write is unnecessary.
This protection is verified indirectly by checking that the smallest input size
still records a non-zero time at every optimization level.

The timed region contains only the authenticated-encryption operation.
For AES-128-GCM the key schedule is performed once before the measurements begin
and is not included.
For Ascon-AEAD128 the single-pass interface performs key setup as part of the
encryption call and cannot separate it, so the full call is timed.
This treats the one-time setup symmetrically: the repeated per-call cost is
measured for both, and the fixed setup that a real deployment would perform once
is excluded from both where the interface allows.
All measurements encrypt plaintext with no associated data.

Each measured iteration is streamed over the serial link as a row of
comma-separated values containing the algorithm, the board, the clock frequency,
the input size, the iteration index, and the measured ticks.
Summary statistics, namely the median, minimum, maximum, and interquartile
range, are computed offline rather than on the device, which keeps the on-device
harness minimal and allows the analysis to be revised without re-running the
measurements.
Under interrupt masking the per-iteration variation was zero on both boards for
all measured workloads, so the median, minimum, and maximum coincide.

\section{Benchmark Parameters}

\subsection{Input Sizes}

Message sizes are swept on a logarithmic scale.
On the L476 the sizes are 1, 10, 100, 1000, and 10\,000 bytes.
On the F072 the largest size is omitted, because the buffers required for a
10\,000-byte message do not fit within 16\,KB of SRAM, so the F072 sizes are
1, 10, 100, and 1000 bytes.
The logarithmic spread separates the fixed per-call cost, which dominates at
small sizes, from the per-byte cost, which dominates at large sizes.

Each size is measured for 100 iterations after the discarded warmup run.

\subsection{Optimization-Level Matrix}

Each algorithm on each board is compiled at five optimization levels:
\texttt{-O0}, \texttt{-O1}, \texttt{-O2}, \texttt{-O3}, and \texttt{-Os}.
This sweep shows how much of the measured performance depends on optimization
rather than on the algorithm itself, and it identifies the level at which each
algorithm performs best, which differs by algorithm and by architecture.
Reporting only a single optimization level would hide both effects.

\section{Footprint Measurement}

Code size is measured as a static property of the compiled binaries using
\texttt{arm-none-eabi-size}, which reports the size of the program text, the
initialized data, and the zero-initialized data sections.
Flash usage is taken as the sum of the text and initialized-data sections.

Because a single binary links both algorithms together with the common platform
and harness code, the size of one combined binary cannot attribute code to one
algorithm or the other.
To obtain a per-algorithm figure, three build variants are produced from the
same source through a compile-time selection:
a baseline variant (platform and harness code, no cryptographic implementation),
an AES variant (adds AES-128-GCM), and an ASCON variant (adds Ascon-AEAD128).
The marginal cost of each algorithm is the difference between its variant and
the baseline.

Footprint is reported primarily at the size-optimization level, because that
is the configuration a memory-constrained deployment would use.
Footprint at the other optimization levels is also recorded, as the code size of
the ASCON reference implementation varies considerably with optimization.

\section{Harness Validation}

Before any cryptographic code was measured, the harness was validated with a
\texttt{memcpy} workload of known behavior, measured at sizes 1, 10, 100, and
1000 bytes for 100 iterations each.
Four checks were defined: the measured time should grow roughly in proportion to
the input size; the per-iteration variation should be at most a few ticks; the
smallest input size should record a non-zero time; and the reported clock
frequency should read 48\,MHz on both boards.

The results of this validation are reported in Section~\ref{sec:harness-val}.

\section{Reproducibility and Known Limitations}

The exact mbedTLS configuration is committed to the repository, and the commit
hashes of mbedTLS and the ASCON reference implementation are
recorded~\cite{mbedTLS,ASCONSoftware}.
The compiler, CMake, Ninja, and CubeCLT versions are recorded, as are the
board hardware revisions and the ST-LINK firmware versions.

Several limitations are disclosed for honesty of interpretation.
The reported times include a small fixed overhead from the timer read pair, the
loop branch, and the \texttt{volatile} store.
At the largest size this overhead is on the order of one percent of the total;
at the smallest size it dominates, so the smallest-size figures should be read
as a measurement of fixed per-call cost rather than of encryption work.
The harness does not subtract this overhead, and raw ticks are reported as
measured.
The AES-GCM measurements use the small GHASH table, as disclosed in
Section~3.4.1.
The ART accelerator is left enabled on the L476, which is the realistic
deployment configuration, and cycle counts are reported so that the comparison
does not depend on its effect.
